\documentclass[10pt,a4paper]{article}
\usepackage[margin=1.7cm]{geometry}
\usepackage[T1]{fontenc}
\usepackage[utf8]{inputenc}
\usepackage{array}
\usepackage{booktabs}
\usepackage{longtable}
\usepackage{tabularx}
\usepackage{placeins}
\usepackage{hyperref}
\usepackage{xcolor}
\usepackage{enumitem}
\usepackage{amssymb}

\setlength{\parindent}{0pt}
\setlength{\parskip}{4pt}
\renewcommand{\arraystretch}{1.18}
\hypersetup{colorlinks=true, linkcolor=blue, urlcolor=blue}

\newcommand{\pending}{\textcolor{red}{pending}}
\newcommand{\done}{\textcolor{teal}{done}}

\title{Group 3 Final Report\\NO2 Air Quality Classification across Austrian Monitoring Stations}
\author{Glent Rembeci, Somayeh Zeraati, Muhammed Tayyab Sajjad, Sina Sadeghi}
\date{June 2026}

\begin{document}
\maketitle

\textbf{Report status.} This report is final and completes the required material for all four roles (A, B, C, and D). All tasks have been implemented, verified, and documented.

\section*{1. Title and Authors}
\textbf{Project title:} NO2 Air Quality Classification across Austrian Monitoring Stations.

\noindent\begin{tabular}{p{0.19\linewidth}p{0.25\linewidth}p{0.36\linewidth}p{0.12\linewidth}}
\toprule
Owner & Name & ORCID & Status \\
\midrule
A & Glent Rembeci & \href{https://orcid.org/0009-0003-3981-8301}{0009-0003-3981-8301} & \done \\
B & Somayeh Zeraati & \href{https://orcid.org/0009-0001-3666-0911}{0009-0001-3666-0911} & \done \\
C & Muhammed Tayyab & \href{https://orcid.org/0009-0000-6631-8482}{0009-0000-6631-8482} & \done \\
D & Sina Sadeghi & \href{https://orcid.org/0009-0008-1996-4666}{0009-0008-1996-4666} & \done \\
\bottomrule
\end{tabular}

\section*{2. Abstract}
This project implements a FAIR data science workflow for classifying nitrogen
dioxide (NO2) air quality levels across Austrian monitoring stations. Hourly 2024
measurements from the European Environment Agency were normalized into DBRepo and
used to train a Random Forest classifier that predicts whether NO2 concentration
is elevated, defined as \(\geq 40\) ug/m3. The project publishes the source
database, software, trained model, generated outputs, final DMP, and metadata
records with persistent identifiers and cross-links.

\section*{3. DBRepo Schema, ER Diagram, Ontology and Unit Mapping}
Owner A prepared the normalized DBRepo database structure and publication
evidence. The database is available through DBRepo with DOI
\href{https://doi.org/10.82556/3zan-dn41}{10.82556/3zan-dn41}. The repository
contains the schema definition in \texttt{docs/schema.sql}, an ER description in
\texttt{docs/er-diagram.md}, DBRepo metadata notes in
\texttt{docs/dbrepo-metadata-plan.md}, and view/API reuse notes in
\texttt{docs/views.sql}.

Owner B prepared the semantic and unit mapping documentation. The semantic mapping
uses SOSA/SSN for observations, OWL-Time for temporal intervals, ChEBI for NO2,
PROV-O for provenance, and QUDT for units. The detailed mapping is documented in
\texttt{docs/semantic-mapping.md} and \texttt{docs/unit-mapping.md}.

Owner C applied the unit mappings to the DBRepo column metadata via the DBRepo
REST API. The \texttt{measurements.value} column was mapped to
\texttt{qudt:MicroGM-PER-M3}, the \texttt{data\_capture} column to
\texttt{qudt:FRACTION}, and the three timestamp columns (\texttt{start\_time},
\texttt{end\_time}, \texttt{result\_time}) to the SI Digital Framework second unit
(\texttt{si:S}). The mappings were verified via the DBRepo API and documented in
\texttt{docs/unit-mapping.md} and the notebook
\texttt{notebooks/t23\_t25\_unit\_mapping\_and\_data\_loading.ipynb}.

Owner D implemented the denormalized SQL views directly in DBRepo via the REST API and refactored the data ingestion pipeline. The training script was updated to fetch data directly from DBRepo's views, handling pagination and errors robustly instead of relying on local CSVs, which ensures reproducibility.

\noindent\begin{tabular}{p{0.25\linewidth}p{0.30\linewidth}p{0.35\linewidth}p{0.05\linewidth}}
\toprule
Attribute & Meaning & Mapping & Owner \\
\midrule
\texttt{Samplingpoint} & EEA monitoring point identifier & \texttt{sosa:FeatureOfInterest} / \texttt{sosa:Platform} & B \\
\texttt{Pollutant} & NO2 pollutant code & \texttt{sosa:observedProperty}, ChEBI 33101 & B \\
\texttt{Start}, \texttt{End} & Hourly measurement interval & \texttt{sosa:phenomenonTime}, OWL-Time beginning/end & B \\
\texttt{Value} & NO2 concentration & \texttt{sosa:hasSimpleResult}, QUDT mass concentration & B \\
\texttt{Unit} & ug/m3 concentration unit & \texttt{unit:MicroGM-PER-M3} & B \\
\texttt{ResultTime} & Time result was produced & \texttt{sosa:resultTime}, \texttt{prov:generatedAtTime} & B \\
\texttt{value} (DBRepo) & NO2 concentration in DBRepo & \texttt{qudt:MicroGM-PER-M3} applied via REST API & C \\
\texttt{data\_capture} & Data capture fraction & \texttt{qudt:FRACTION} applied via REST API & C \\
\texttt{start\_time}, \texttt{end\_time}, \texttt{result\_time} & Timestamps & SI Digital Framework \texttt{si:S} applied via REST API & C \\
\texttt{v\_measurements\_enriched} & Denormalized view & DBRepo REST API view creation & D \\
\texttt{v\_weekday\_measurements} & Weekday filter view & DBRepo REST API view creation & D \\
\bottomrule
\end{tabular}

\section*{4. Standards Application and Overlap Analysis}
Owner A prepared the RO-Crate and Model Card evidence. The
\texttt{ro-crate-metadata.json} file describes the main software, data, model,
generated outputs, and related identifiers, and the local validation note reports
that the RO-Crate 1.1 required checks passed. The Model Card in
\texttt{docs/model-card.md} documents the classifier purpose, inputs, target
definition, evaluation metrics, ethical limitations, and reuse information.

Owner B prepared the CodeMeta and licensing evidence. The \texttt{codemeta.json}
record documents the software package, authors, license, repository, runtime, and
dependencies. The licensing notes distinguish the EEA input data terms, MIT
software license, CC BY 4.0 generated outputs/model documentation, and reuse
conditions for published artefacts.

Owner C prepared the standards overlap analysis in
\texttt{docs/standards-overlap-analysis.md}. The analysis covers all ten pairwise
comparisons between RO-Crate, CodeMeta, FAIR4ML, Croissant, and Model Card,
identifying shared fields, unique fields, and conflicts. Key findings: RO-Crate
acts as the integration layer referencing all other standards; FAIR4ML and Model
Card have the highest field overlap; CodeMeta and Croissant are the most
complementary pair; and licence management is a cross-cutting concern requiring
careful alignment across all five standards.

Owner C also prepared the FAIR4ML metadata file in
\texttt{docs/fair4ml-metadata.json}, documenting the Random Forest classifier
algorithm, all seven input features, hyperparameters, evaluation metrics
(accuracy 0.8005, precision 0.2107, recall 0.7251, F1 0.3266), intended use, and
known limitations.

Owner D prepared the Croissant metadata in \texttt{croissant.json}. The record
documents the dataset structure and field mappings (Value, Start, End,
Samplingpoint) to the raw EEA Parquet source files, including the unit-of-measurement
references for the measured values. Owner D also managed the GitHub--Zenodo integration and the \texttt{CITATION.cff} setup for versioned software releases.

\section*{5. DMP Comparison}

Owner B prepared the initial-vs-final DMP comparison (T4.5). The initial DMP
(prepared in Part~2) is available at
\href{https://handle.test.datacite.org/10.70124/3gj7h-t2832}{10.70124/3gj7h-t2832}.
The final DMP record is available at
\href{https://doi.org/10.70124/pcepy-07563}{10.70124/pcepy-07563}.
The full comparison is documented in \texttt{docs/dmp-comparison-template.md}.

\textbf{1. What the initial DMP promised that turned out to be unrealistic.}

The main thing we underestimated in the initial DMP was how much coordination
persistent identifiers would require. We wrote that data would be published
openly with DOIs, but the initial plan assumed this would be one or two steps.
In practice, we ended up managing five separate artefact records: the DBRepo
database, the GitHub/Zenodo software release, the TUWRD model deposit, the
TUWRD generated-output deposit, and the final DMP record -- each needing its
own metadata, licence statement, and cross-references to the others.

Licensing was also more work than the initial plan suggested. We had to check
the exact EEA source data reuse conditions, choose a compatible MIT licence for
the code, and separately decide on CC~BY~4.0 for the outputs, including checking
whether ShareAlike clauses from the source data would apply. The initial DMP did
not mention DBRepo at all, so the schema design, semantic mappings, unit
mappings via the REST API, and the full reimplementation of the data pipeline
from local CSV files to DBRepo views were all absent from the initial plan.

\textbf{2. What writing the final DMP actually required that the initial one glossed over.}

The biggest gap was that the final DMP needed to show what we actually did,
not just what we planned to do. For us that meant producing five metadata
artefact files (RO-Crate, CodeMeta, Croissant, FAIR4ML, and a Model Card)
and making sure they were consistent with each other. The initial DMP did
not mention any of these standards.

We also had to cross-link every persistent identifier across all records.
This meant putting the DBRepo DOI
(\href{https://doi.org/10.82556/3zan-dn41}{10.82556/3zan-dn41}),
the Zenodo DOI
(\href{https://doi.org/10.5281/zenodo.20432795}{10.5281/zenodo.20432795}),
the TUWRD model DOI
(\href{https://doi.org/10.70124/ye7mg-p5v03}{10.70124/ye7mg-p5v03}),
the TUWRD generated-data DOI
(\href{https://doi.org/10.70124/4jwcg-70c34}{10.70124/4jwcg-70c34}),
and the DMP DOI itself
(\href{https://doi.org/10.70124/pcepy-07563}{10.70124/pcepy-07563})
into the maDMP JSON, the TUWRD deposit records, and the RO-Crate. Keeping
these in sync across five different systems took noticeably more time than
we expected when writing the initial plan.

Owner C wrote the final DMP in DAMAP following the FWF template (T4.1). The DMP
describes all three dataset groups, storage locations (GitHub, Zenodo, TUWRD,
DBRepo), access rights, licences, and reuse planning. The DMP was exported and
submitted as part of the final TUWRD DMP record.

Owner D reviewed the final DMP to ensure alignment with the deployed FAIR
implementations (T4.2), verifying that the stated metadata standards exactly
matched the implemented repository capabilities and documentation.

\section*{6. OSTrails Test Results Table}
The group assessment service was implemented under
\texttt{groups/group-3} in the course GitLab repository. The merge request is
\href{https://gitlab.tuwien.ac.at/crdm/data-stewardship/dast-2026-ue/-/merge_requests/15}{GitLab MR 15}.

\noindent\begin{tabular}{p{0.14\linewidth}p{0.14\linewidth}p{0.12\linewidth}p{0.44\linewidth}}
\toprule
Metric ID & Owner & Outcome & Message \\
\midrule
T-DCSC-026 & A & pass & Implemented and checked with local pass/fail examples for the assigned DMP evaluation test. \\
T-DCSC-037 & A & pass & Implemented and checked with local pass/fail examples for the assigned DMP evaluation test. \\
T-DCSC-065 & A & pass & Implemented and checked with local pass/fail examples for the assigned DMP evaluation test. \\
T-DCSC-058 & B & pass & Implemented and checked with local pass/fail examples for the assigned DMP evaluation test. \\
T-DCSC-069 & B & pass & Implemented and checked with local pass/fail examples for the assigned DMP evaluation test. \\
T-DCSC-083 & B & pass & Implemented and checked with local pass/fail examples for the assigned DMP evaluation test. \\
T-DCSC-001 & C & pass & Field `is\_reused' found on at least one dataset. GitLab CI pipeline passed. \\
T-DCSC-016 & C & pass & At least one dataset entry has no is\_reused field and is treated as new data. GitLab CI pipeline passed. \\
T-DCSC-032 & C & pass & The maDMP JSON document validates against the RDA DMP Common Standard Schema (v1.1). GitLab CI pipeline passed. \\
T-DCSC-003 & D & pass & Implemented and checked with local pass/fail examples for the assigned DMP evaluation test. \\
T-DCSC-022 & D & pass & Implemented and checked with local pass/fail examples for the assigned DMP evaluation test. \\
T-DCSC-090 & D & pass & Implemented and checked with local pass/fail examples for the assigned DMP evaluation test. \\
\bottomrule
\end{tabular}

\section*{7. Individual Contribution and Sign-off Tables}
\subsection*{Owner A: Glent Rembeci --- Matriculation No.: e12332209}
\noindent\begin{tabular}{p{0.08\linewidth}p{0.42\linewidth}p{0.28\linewidth}p{0.1\linewidth}}
\toprule
Task & What I did & Link / DOI / commit & Done? \\
\midrule
T2.1 & Designed the normalized relational schema for the EEA NO2 data in 3NF. Produced SQL CREATE statements, an ER diagram, and descriptive database metadata. Created and configured the DBRepo database entry with publisher, licence, and provenance metadata. & \href{https://doi.org/10.82556/3zan-dn41}{DOI 10.82556/3zan-dn41}, \href{https://github.com/tayyabsajjad3/no2-air-quality-classification-austria/blob/main/docs/schema.sql}{docs/schema.sql} & \checkmark \\
T3.1 & Created the \texttt{ro-crate-metadata.json} file describing the full experiment package: input datasets, code, trained model, generated outputs, authors, licences, and relationships. Validated with \texttt{rocrate-validator 0.9.0}; all 38 required RO-Crate 1.1 checks passed. & \href{https://github.com/tayyabsajjad3/no2-air-quality-classification-austria/blob/main/ro-crate-metadata.json}{ro-crate-metadata.json}, \href{https://github.com/tayyabsajjad3/no2-air-quality-classification-austria/blob/main/docs/validation/}{docs/validation/} & \checkmark \\
T3.5 & Wrote the Model Card at \texttt{docs/model-card.md} covering model description, intended use, out-of-scope uses, training data with DOI, evaluation results with a precision/recall/F1 table, limitations, ethical considerations, and licence. All required sections are at least three sentences. & \href{https://github.com/tayyabsajjad3/no2-air-quality-classification-austria/blob/main/docs/model-card.md}{docs/model-card.md} & \checkmark \\
T3.9 & Uploaded the trained model artefact (\texttt{no2\_air\_quality\_classifier.joblib}) and model-card to the TUWRD test instance as a separate model-only deposit. Set resource type to Model, added all four creators with ORCIDs, CC~BY~4.0 licence, and cross-links to Zenodo, GitHub, DBRepo, and the generated-data deposit. & \href{https://doi.org/10.70124/ye7mg-p5v03}{DOI 10.70124/ye7mg-p5v03} & \checkmark \\
T4.4 & Submitted the final DMP PDF and maDMP JSON as a single record in the DaSt-2026-final community on TUWRD. Set resource type to Data Management Plan, added all four creators with ORCIDs, CC~BY~4.0 licence, and related identifiers pointing to GitHub, Zenodo, DBRepo, and both TUWRD deposits. & \href{https://doi.org/10.70124/pcepy-07563}{DOI 10.70124/pcepy-07563} & \checkmark \\
T5.1 & Created the group-3 GitLab merge request for the shared OSTrails assessment service, incorporating all four members' tests. Ensured a single pull request per the exercise instructions. & \href{https://gitlab.tuwien.ac.at/crdm/data-stewardship/dast-2026-ue/-/merge_requests/15}{GitLab MR 15} & \checkmark \\
T5.2 & Implemented OSTrails DMP evaluation tests T-DCSC-026, T-DCSC-037, and T-DCSC-065. Verified correct pass/fail outcomes using local example maDMP files before submission. & \href{https://gitlab.tuwien.ac.at/crdm/data-stewardship/dast-2026-ue/-/merge_requests/15}{GitLab MR 15} & \checkmark \\
\bottomrule
\end{tabular}

\textbf{Sign-off chain:} Owner A confirmed by Owner B --- \done

\newpage
\subsection*{Owner B: Somayeh Zeraati --- Matriculation No.: e12353396}
\noindent\begin{tabular}{p{0.08\linewidth}p{0.42\linewidth}p{0.28\linewidth}p{0.1\linewidth}}
\toprule
Task & What I did & Link / DOI / commit & Done? \\
\midrule
T2.2 & Mapped all dataset attributes to semantic concepts in domain-specific ontologies (SOSA/SSN, OWL-Time, ChEBI, PROV-O, QUDT). Justified ontology choices in 2--3 sentences each. Added the semantic mappings to the DBRepo column metadata via the REST API. & \href{https://github.com/tayyabsajjad3/no2-air-quality-classification-austria/blob/main/docs/semantic-mapping.md}{docs/semantic-mapping.md} & \checkmark \\
T2.7 & Prepared the second GitHub release (v2.0) incorporating all WP2 changes. Wrote descriptive release notes covering schema, semantic mappings, unit mappings, DBRepo data loading, API reimplementation, and SQL views. & \href{https://github.com/tayyabsajjad3/no2-air-quality-classification-austria/releases}{GitHub releases} & \checkmark \\
T3.2 & Created \texttt{codemeta.json} following the CodeMeta~2.0 schema. Included name, version, all four authors with ORCIDs, MIT licence SPDX identifier, Python runtime, all pinned dependencies from \texttt{requirements.txt}, and the GitHub repository URL. & \href{https://github.com/tayyabsajjad3/no2-air-quality-classification-austria/blob/main/codemeta.json}{codemeta.json} & \checkmark \\
T3.6 & Assigned licences to all three artefact categories. Verified EEA input data reuse terms (CC~BY~4.0), confirmed compatibility with the MIT software licence, and selected CC~BY~4.0 for generated outputs and the trained model. Documented all decisions with rationale in \texttt{docs/licences.md}. & \href{https://github.com/tayyabsajjad3/no2-air-quality-classification-austria/blob/main/docs/licences.md}{docs/licences.md} & \checkmark \\
T3.10 & Uploaded the generated output data (predictions, evaluation figures, confusion matrix, metrics) to the TUWRD test instance as a separate dataset deposit. Set resource type to Dataset, added all four creators with ORCIDs, CC~BY~4.0 licence, and cross-links to Zenodo, GitHub, DBRepo, and the model deposit. & \href{https://doi.org/10.70124/4jwcg-70c34}{DOI 10.70124/4jwcg-70c34} & \checkmark \\
T4.3 & Exported the machine-actionable DMP from the TU Wien DMP Tool as a JSON document. Added missing identifiers (DBRepo DOI, Zenodo DOI, TUWRD deposit DOIs, initial DMP DOI) in accordance with the RDA maDMP recommendation. Included the maDMP JSON in the final TUWRD DMP record. & \href{https://doi.org/10.70124/pcepy-07563}{DOI 10.70124/pcepy-07563} & \checkmark \\
T4.5 & Compared the initial DMP from Part~2 with the final DMP. Answered the two required questions: what the initial plan overpromised and what the final plan actually required in effort, decisions, and knowledge. Both DMP links are referenced in Section~5 of this report. & \href{https://github.com/tayyabsajjad3/no2-air-quality-classification-austria/blob/main/docs/dmp-comparison-template.md}{docs/dmp-comparison-template.md} & \checkmark \\
T5.3 & Implemented OSTrails DMP evaluation tests T-DCSC-058, T-DCSC-069, and T-DCSC-083. Verified correct pass/fail outcomes using local example maDMP files before submission. & \href{https://gitlab.tuwien.ac.at/crdm/data-stewardship/dast-2026-ue/-/merge_requests/15}{GitLab MR 15} & \checkmark \\
\bottomrule
\end{tabular}

\textbf{Sign-off chain:} Owner B confirmed by Owner C --- \done

\newpage
\subsection*{Owner C: Muhammed Tayyab Sajjad --- Matriculation No.: e12510765}

\noindent\begin{tabular}{p{0.08\linewidth}p{0.42\linewidth}p{0.28\linewidth}p{0.1\linewidth}}
\toprule
Task & What I did & Link / DOI / commit & Done? \\
\midrule
T2.3 & Applied unit-of-measurement ontology mappings to five DBRepo columns via the REST API: value to QUDT MicroGM-PER-M3, data\_capture to QUDT FRACTION, and start\_time/end\_time/result\_time to SI Digital Framework second. Verified mappings via the DBRepo API. & \href{https://github.com/tayyabsajjad3/no2-air-quality-classification-austria/blob/main/docs/unit-mapping.md}{docs/unit-mapping.md} & \checkmark \\
T2.5 & Loaded all raw EEA Parquet data into DBRepo using the Python dbrepo client. All 8 tables were populated: measurements (140,160 rows), sampling\_points (16), and all lookup tables. Verified data via the DBRepo API and web UI. & \href{https://github.com/tayyabsajjad3/no2-air-quality-classification-austria/blob/main/notebooks/t23_t25_unit_mapping_and_data_loading.ipynb}{notebooks/t23\_t25} & \checkmark \\
T3.3 & Created FAIR4ML metadata file documenting the Random Forest classifier: algorithm, 7 input features, hyperparameters, evaluation metrics (accuracy 0.8005, precision 0.2107, recall 0.7251, F1 0.3266), intended use, and limitations. & \href{https://github.com/tayyabsajjad3/no2-air-quality-classification-austria/blob/main/docs/fair4ml-metadata.json}{docs/fair4ml-metadata.json} & \checkmark \\
T3.7 & Rewrote the README to include abstract, contributors table with ORCIDs, all three licences, requirements and installation instructions, step-by-step reproduction instructions for both local and DBRepo execution, and inputs/outputs tables. Resolved two merge conflicts. & \href{https://github.com/tayyabsajjad3/no2-air-quality-classification-austria/blob/main/README.md}{README.md} & \checkmark \\
T3.8 & Created \texttt{CITATION.cff} in the repository root with all four contributors, ORCIDs, and Zenodo DOI. Set up the Zenodo--GitHub integration for automatic archiving on release. Owner D added the Zenodo DOI badge to the README. & \href{https://github.com/tayyabsajjad3/no2-air-quality-classification-austria/blob/main/CITATION.cff}{CITATION.cff} & \checkmark \\
T3.11 & Produced standards overlap analysis covering all ten pairwise comparisons between RO-Crate, CodeMeta, FAIR4ML, Croissant, and Model Card. Identified shared fields, unique fields, conflicts, and key findings. & \href{https://github.com/tayyabsajjad3/no2-air-quality-classification-austria/blob/main/docs/standards-overlap-analysis.md}{docs/standards-overlap-analysis.md} & \checkmark \\
T4.1 & Wrote the final DMP in DAMAP following the FWF template. Added all four contributors, three dataset groups, four storage locations, access rights, licences, and reuse planning. Exported the DMP using the FWF template. & \href{https://dmptool.tuwien.ac.at/dmp/13011}{dmptool.tuwien.ac.at/dmp/13011} & \checkmark \\
T5.4 & Implemented three OSTrails DMP evaluation tests (T-DCSC-001, T-DCSC-016, T-DCSC-032) as a Node.js/Express REST API with Docker packaging. All three tests passed the GitLab CI/CD pipeline. & \href{https://gitlab.tuwien.ac.at/crdm/data-stewardship/dast-2026-ue/-/merge_requests/15}{GitLab MR 15} & \checkmark \\
T6.1 & Reviewed Group 5's DBRepo database: checked 3 tables, 3 views, and semantic annotations. Found no unit URIs or concept mappings on any column. Contributed section (b) of the group review report. & \href{https://github.com/tayyabsajjad3/no2-air-quality-classification-austria/tree/main/DMP\%20review}{DMP review/} & \checkmark \\
\bottomrule
\end{tabular}

\textbf{Sign-off chain:} Owner C confirmed by Owner D --- \done

\newpage
\subsection*{Owner D: Sina Sadeghi --- Matriculation No.: e12507323}

\noindent\begin{tabular}{p{0.08\linewidth}p{0.42\linewidth}p{0.28\linewidth}p{0.1\linewidth}}
\toprule
Task & What I did & Link / DOI / commit & Done? \\
\midrule
T1.4 & Made the first GitHub release of the project using the GitHub release functionality. Set up the repository structure, Python environment with \texttt{uv}, and ensured the initial codebase was clean and reproducible before tagging the first release. & Commit \texttt{02a3027} & \checkmark \\
T2.4 & Implemented and created the denormalized SQL views directly in DBRepo via the REST API for downstream ML use. & Commit \texttt{b7d8852} & \checkmark \\
T2.6 & Refactored the data ingestion and ML pipeline to fetch data directly from DBRepo's REST API views (with pagination and error handling) instead of local CSVs. & Commit \texttt{4d1118e} & \checkmark \\
T3.4 & Created the Croissant metadata file (\texttt{croissant.json}) documenting the dataset structure and semantic field mappings to the raw EEA Parquet source files. & \href{https://github.com/tayyabsajjad3/no2-air-quality-classification-austria/blob/main/croissant.json}{croissant.json} & \checkmark \\
T3.8 & Added the Zenodo DOI badge to the README and verified the GitHub--Zenodo integration was correctly linked. Supported Owner C in setting up \texttt{CITATION.cff} and confirmed that the release triggered automatic Zenodo archiving. & \href{https://github.com/tayyabsajjad3/no2-air-quality-classification-austria/blob/main/CITATION.cff}{CITATION.cff} & \checkmark \\
T4.2 & Reviewed and aligned the final DMP implementations, verifying that DBRepo capabilities and metadata documentation matched the submitted FWF template. & \href{https://dmptool.tuwien.ac.at/dmp/13011}{dmptool.tuwien.ac.at/dmp/13011} & \checkmark \\
T5.5 & Implemented three OSTrails DMP evaluation tests (T-DCSC-003, T-DCSC-022, T-DCSC-090). Verified locally and via GitLab CI/CD pipeline. & \href{https://gitlab.tuwien.ac.at/crdm/data-stewardship/dast-2026-ue/-/merge_requests/15}{GitLab MR 15} & \checkmark \\
T6.1 & Reviewed Group 5's metadata standard artefacts (RO-Crate, CodeMeta, FAIR4ML, Croissant, Model Card), noting high-quality documentation but DBRepo UI inconsistencies. & \href{https://github.com/tayyabsajjad3/no2-air-quality-classification-austria/tree/main/DMP\%20review}{DMP review/} & \checkmark \\
T6.2 & Assembled and formatted the final review report for Group 5 in LaTeX, harmonized the ratings across all checks, and prepared the final PDF. & \href{https://github.com/tayyabsajjad3/no2-air-quality-classification-austria/tree/main/DMP\%20review/report}{DMP review/report} & \checkmark \\
T7.1 & Assembled and formatted the final Group 3 project report in LaTeX, harmonized statuses, and ensured all tasks were correctly assigned and documented. & \href{https://github.com/tayyabsajjad3/no2-air-quality-classification-austria/tree/main/docs/final-report}{docs/final-report} & \checkmark \\
\bottomrule
\end{tabular}

\textbf{Sign-off chain:} Owner D confirmed by Owner A --- \done

\end{document}